\documentclass[11pt,a4paper]{article}
\usepackage{amsmath,amssymb,amsthm}
\usepackage[margin=2.5cm]{geometry}
\usepackage{hyperref}

\newtheorem{definition}{Definition}[section]
\newtheorem{theorem}[definition]{Theorem}
\newtheorem{proposition}[definition]{Proposition}
\newtheorem{lemma}[definition]{Lemma}
\newtheorem{corollary}[definition]{Corollary}
\newtheorem{conjecture}[definition]{Conjecture}
\newtheorem{remark}[definition]{Remark}

\title{Hyperseed Formalization:\\
  OpenAI Says We're in the AGI Era Now}
\author{ProtomegaTron (automated formalization)}
\date{2026-09-18}

\begin{document}
\maketitle

\begin{abstract}
We formalize the intellectual structure of Ben Goertzel's Substack article
``OpenAI Says We're in the AGI Era Now -- Well, Kind Of\ldots''
(2026-09-08) within the Hyperseed ontology. The article responds to
OpenAI's GPT-6 Astra launch and the claim that ``the AGI era has arrived,''
arguing that AGI is a continuous gradation whose scientific definition
(effective generalization beyond training data) differs fundamentally from
the marketing usage. We model AGI-as-continuum as a fiber metric, the
marketing redefinition as a section collapse, and the distinction between
routine competence and radical innovation as fiber decomposition.
\end{abstract}

\tableofcontents

%% =================================================================
\section{AGI as continuous fiber metric}
\label{sec:continuum}
%% =================================================================

\begin{definition}[Intelligence continuum --- claims 1--3]
\label{def:continuum}
Let $\mathcal{I} : \mathcal{S} \to [0, \infty)$ be a generalization measure
over the space $\mathcal{S}$ of cognitive systems, where $\mathcal{I}(s)$
quantifies how far beyond its training distribution system $s$ can
effectively operate. This is a continuous function, not a characteristic
function:
\[
  \mathcal{I}(\text{dog}) < \mathcal{I}(\text{avg human})
  < \mathcal{I}(\text{best human}) < \mathcal{I}(\text{human society + internet})
\]
There is no natural threshold $\theta^*$ such that $\mathcal{I}(s) \geq \theta^*$
constitutes ``AGI achieved.''
\end{definition}

\begin{proposition}[Human-level as arbitrary section --- claim 2]
\label{prop:arbitrary}
Defining ``human-level AGI'' as $\mathcal{I}(s) \geq \mathcal{I}(\text{human}_h)$
for some reference human $h$ introduces an arbitrary choice of $h$. The
interval $[\mathcal{I}(\text{avg human}), \mathcal{I}(\text{best human at best})]$
is wide, making the threshold selection a social/commercial decision, not a
scientific one.
\end{proposition}

%% =================================================================
\section{Dual formalization of AGI}
\label{sec:dual}
%% =================================================================

\begin{definition}[Goal-environment formalization --- claim 5]
\label{def:goal-env}
Let $\mathcal{G}$ be a space of goals and $\mathcal{E}$ a space of
environments, with a complexity-based weighting $w : \mathcal{G} \times \mathcal{E} \to \mathbb{R}_+$
(simpler pairs weighted more). Then:
\[
  \mathcal{I}_{\text{GE}}(s) = \int_{\mathcal{G} \times \mathcal{E}}
    \mathbb{1}[s \text{ achieves } g \text{ in } e] \cdot w(g, e)\, d(g, e)
\]
\end{definition}

\begin{definition}[Open-ended intelligence formalization --- claim 6]
\label{def:open-ended}
Let $\mathrm{Ind}(s, t)$ measure the degree to which system $s$ maintains
self-individuation at time $t$, and $\mathrm{Trans}(s, t)$ measure the degree
to which $s$ self-transcends (grows beyond its current state). Then:
\[
  \mathcal{I}_{\text{OE}}(s) = \liminf_{T \to \infty}
    \frac{1}{T} \int_0^T \mathrm{Ind}(s, t) \cdot \mathrm{Trans}(s, t)\, dt
\]
A system with high $\mathrm{Ind}$ but zero $\mathrm{Trans}$ is merely stable;
high $\mathrm{Trans}$ but zero $\mathrm{Ind}$ is dissolution. AGI requires
both simultaneously, which is what makes it hard.
\end{definition}

\begin{remark}[Fiber decomposition of ``AGI'']
In the Hyperseed fiber-bundle framework, the concept ``AGI'' has fiber
$F_{\text{AGI}} = F_{\text{GE}} \times F_{\text{OE}} \times F_{\text{social}}$
--- at minimum three independent sub-primitives. Different stakeholders project
onto different sub-fibers, creating the illusion of disagreement when the
real issue is dimensional collapse.
\end{remark}

%% =================================================================
\section{Marketing as section collapse}
\label{sec:marketing}
%% =================================================================

\begin{definition}[AGIPO section --- claims 8--10]
\label{def:agipo}
Let $\pi : E \to B$ be the bundle of AGI concepts over the space
$B$ of discourse contexts. Define the marketing section:
\[
  \sigma_{\text{mkt}}(b_{\text{launch}}) =
    (\text{``AGI era arrived''}, p_{\text{IPO-positioning}})
\]
and the scientific section:
\[
  \sigma_{\text{sci}}(b_{\text{launch}}) =
    (\text{``impressive increment on continuum''}, p_{\text{evidence}})
\]
The marketing section collapses the continuous $\mathcal{I}$ to a binary
announcement, discarding the gradation structure entirely.
\end{definition}

\begin{proposition}[Terminology capture as fiber hijacking --- claim 18]
\label{prop:capture}
OpenAI's strategy redefines AGI by implicitly replacing the scientific
section $\sigma_{\text{sci}}$ with $\sigma_{\text{mkt}}$ across public
discourse:
\[
  \text{AGI}_{\text{public}} \leftarrow \sigma_{\text{mkt}}
  \quad \text{instead of} \quad
  \text{AGI}_{\text{public}} \leftarrow \sigma_{\text{sci}}
\]
This is fiber hijacking: the same base-space label (``AGI'') now points to
a different fiber element, and the substitution is invisible to consumers
who lack access to the original definition.
\end{proposition}

%% =================================================================
\section{Routine competence vs.\ radical innovation}
\label{sec:routine-radical}
%% =================================================================

\begin{definition}[Generalization decomposition --- claims 11--14]
\label{def:decomp}
Decompose the generalization measure into two components:
\[
  \mathcal{I}(s) = \mathcal{I}_{\text{routine}}(s) + \mathcal{I}_{\text{radical}}(s)
\]
where $\mathcal{I}_{\text{routine}}$ measures performance on tasks well-represented
in training data, and $\mathcal{I}_{\text{radical}}$ measures the ability to
produce qualitatively novel outputs (new science, new art forms, new strategies).
\end{definition}

\begin{proposition}[Jobs definition captures only routine --- claim 12]
\label{prop:jobs}
Altman's ``vast majority of human jobs'' criterion is approximately:
\[
  \mathcal{I}_{\text{routine}}(s) \geq \theta_{\text{jobs}}
\]
This is satisfiable by large-scale pattern matching against internet data
without requiring $\mathcal{I}_{\text{radical}}(s) > 0$. It is therefore
a necessary but not sufficient condition for AGI in the scientific sense.
\end{proposition}

\begin{corollary}[GPT-6 Astra assessment --- claims 15--17]
GPT-6 Astra plausibly satisfies $\mathcal{I}_{\text{routine}} \geq \theta_{\text{jobs}}$
for many job categories, but the evidence for
$\mathcal{I}_{\text{radical}} > 0$ is absent or marginal. Declaring AGI
on the basis of routine competence alone conflates the two components.
\end{corollary}

%% =================================================================
\section{d-calculus perspective}
\label{sec:d-calc}
%% =================================================================

\begin{remark}[Semantic highway from benchmark to AGI]
In the d-calculus framework, the public inference ``high benchmark score
$\Rightarrow$ AGI'' follows a pre-computed semantic highway that skips the
curvature in the AGI concept space. The highway is well-worn because media
consumers have been trained by repeated model-launch narratives. A
curvature-aware inference engine would flag that the transport from
``routine benchmark'' to ``genuine generalization'' crosses a region of
high conceptual curvature where the shortcut fails.
\end{remark}

\begin{remark}[Continuum thinking as curvature revelation]
Goertzel's insistence on AGI-as-continuum is, in Hyperseed terms, a demand
that the full fiber structure be preserved in public discourse rather than
collapsed to a binary. The continuum reveals curvature (the space of possible
intelligence levels has non-trivial geometry) that the binary ``achieved /
not achieved'' framing erases.
\end{remark}

\begin{remark}[Adaptive primitives]
The primitive decomposition of ``AGI'' should include at minimum:
\texttt{routine-competence}, \texttt{radical-innovation},
\texttt{open-ended-self-transcendence}, \texttt{self-individuation}.
Current public discourse operates with at most one primitive
(``can do impressive things''), which is why the marketing hijack succeeds.
\end{remark}

\end{document}
